\section{Introduction}
\label{sec:intro}

\textit{Software compartmentalization} hardens application security by enforcing the \textit{principle of least privilege}~\cite{protection-info-system, sok-compartmentalization}. A compartmentalized application divides its code and data into \textit{compartments} and isolates each one in its own protection domain, using mechanisms that range from process isolation to WebAssembly sandboxing and hardware memory tagging~\cite{provos2003preventing, rlbox, isolation-without-taxation, xfi, bgi, erim, hacks, bulkhead}. \rev{Compartmentalization is increasingly retrofitted into existing software, such as browser libraries~\cite{rlbox} and kernel drivers~\cite{sud, nooks, microdrivers, ksplit, lvds, lxds}, to confine the damage that a compromised component can cause. Where it has been deployed, it has reduced the impact of real vulnerabilities, for example in the sandbox compartments of OpenSSH and slaacd in OpenBSD~\cite{sok-compartmentalization}.}

\rev{However, isolation alone does not make a compartmentalized application
secure. Compartments still communicate through a \textit{compartment
interface}, the functions that one compartment invokes in another and the data
that these functions exchange~\cite{sok-compartmentalization, assess_civ}.
Isolation prevents a compartment from accessing another compartment's memory
directly, but it does not control this data. A compromised compartment can
therefore attack other compartments without escaping its isolation, by abusing
the data that the interface legitimately exposes. Prior work has found such
flaws in real software. In Firefox, a sandboxed image library could read and
write arbitrary memory of the renderer if the renderer did not check callback
arguments~\cite{rlbox}. Similarly, a sandboxed libghostscript can write to
arbitrary memory of ImageMagick, and an attacker can extract the keys of an
isolated OpenSSL through a single interface flaw~\cite{assess_civ}. These
attacks do not need to hijack control flow. Corrupted data alone, such as a
pointer that the application later dereferences, is enough~\cite{dui}.}

\rev{Such flaws are called \textit{compartment interface vulnerabilities}
(\vulName{}s)~\cite{assess_civ}. They are instances of the confused-deputy
problem, in which an untrusted compartment manipulates the interface so that a
trusted compartment misuses its authority on the attacker's
behalf~\cite{hardy1988confuseddeputy}. \vulName{}s are also common. Fuzzing 39
compartmentalization scenarios across 25 applications revealed 629 potential
\vulName{}s~\cite{assess_civ}, and a study of 11 Linux drivers under isolation
finds several classes of \vulName{}s prevalent across these
drivers~\cite{huang_sok_attack_surface_device_drv_iso}.}

\rev{This paper addresses \textit{insecure data-flow} \vulName{}s, which leak or
corrupt interface data and account for most of the \vulName{}s that these
studies report~\cite{huang_sok_attack_surface_device_drv_iso, assess_civ}.
\vulName{}s that abuse the order or timing of interface calls are outside our
scope. \textit{Data leakage} occurs when a trusted compartment exposes more
interface data than an untrusted compartment needs. \textit{Data corruption}
occurs when an untrusted compartment modifies interface data that a trusted one
later uses. Preventing both requires three decisions for every piece of data
that crosses an interface. The first is which data a trusted compartment may
pass to an untrusted one. The second is which updates by the untrusted
compartment may reach the trusted state. The third is which pointers from the
untrusted compartment must be validated before the trusted compartment
dereferences them. Real interfaces make these decisions hard to get right by
hand. For example, the 214 OpenSSL functions that Nginx calls exchange more
than 350,000 interface fields (Section~\ref{subsec:eval_policy_coverage}).}

\rev{Prior work has measured this attack surface and built tools that detect
\vulName{}s~\cite{assess_civ, civscope, dui,
huang_sok_attack_surface_device_drv_iso}, but protecting interfaces remains
largely manual. Almost all compartmentalization frameworks neglect the security
of interfaces~\cite{assess_civ}. Even RLBox, the framework that sandboxes
libraries in Firefox, relies extensively on validation code that developers
write by hand~\cite{assess_civ}, and its authors found writing this code the
hardest part of retrofitting Firefox~\cite{rlbox}. Driver isolation systems
that copy only part of the shared state choose that part for compatibility and
performance, not to prevent \vulName{}s~\cite{microdrivers, ksplit}. Where
interface checks are automated, they rely on invariants inferred from observed
executions~\cite{microdriver+} or on rules that developers
write~\cite{Mao11LXFI}, and inferring sound and complete invariants remains
challenging~\cite{huang_sok_attack_surface_device_drv_iso}. \textbf{As a
result, the data that crosses a compartment interface remains a largely
unprotected attack surface, and protecting it is left to the developer.}}

\rev{We present \sysName{}, which makes the three decisions automatically and
enforces them at runtime. Our key observation is that the code on the two
sides of an interface already shows how each field of the interface data is
used. A field that the untrusted code never reads, or the trusted code never
writes, need not be passed to the untrusted compartment. An update to a field
that the untrusted code never writes, or the trusted code never reads, need not
be copied back. A pointer that the untrusted code may modify and the trusted
code later dereferences must be validated first. \sysName{} computes these
access patterns with static analysis on a Program Dependence Graph
(PDG)~\cite{Ferrante87PDG, LiuTJ17Ptrsplit} and derives one policy for each
decision. A \emph{least-privilege read} policy (SP1) keeps data that the
untrusted compartment does not need out of it, providing confidentiality. A
\emph{least-privilege write} policy (SP2) discards updates that the trusted
compartment should not accept, providing integrity. A \emph{pointer integrity}
policy (SP3) checks that each pointer from the untrusted compartment refers to
a live object of the expected type in the expected compartment's memory and
that the access stays within the object's bounds.}

\rev{\sysName{} enforces these policies by generating wrappers around interface
calls, integrated with RLBox~\cite{rlbox}. The wrappers copy only the data
that SP1 and SP2 allow, and they validate pointers against a trusted object
registry that records the type, bounds, and lifetime of known objects in both
compartments' memory.} Our key contributions are listed below.

%\sysname{}'s analysis automatically computes a set of interface fields that the trusted compartment does not need to expose to the untrusted one, providing confidentiality for those fields. It also computes a set of interface fields that 
% GT: removed the following as it's too detailed in the intro.
%If the trusted compartment writes a field and the untrusted one reads it, that field is treated as data passed from T to U. If untrusted code writes a field and trusted code reads it, that field is treated as data passed from U to T. If U may modify a pointer value that trusted code later dereferences, trusted code must validate that the pointer references the expected memory domain before use.
%Instead of copying whole interface objects into the sandbox, the wrapper passes only fields that U reads. Instead of propagating whole objects back to T after a sandbox call, it propagates only fields whose values T reads. For pointer fields whose values U may modify and T later dereferences, \sysName{} inserts runtime domain checks. 

\begin{itemize}
    \item We analyze insecure data-flow \vulName{}s and
    define three security policies (SP1, SP2, SP3) that address data
    leakage, data corruption, and \rev{pointer integrity}
    (Section~\ref{sec:CIV_taxonomy_mitigation}).

    \item We design and implement a static analysis framework that
    infers these policies from a compartment boundary specification by
    analyzing how trusted and untrusted compartments access interface data (Section~\ref{sec:impl}).

    \item We integrate the inferred policies with RLBox by generating
    boundary wrappers that restrict data sharing at field granularity and validate \rev{pointer values against a trusted object registry} (Section~\ref{sec:impl}).

    \item We validate the enforced policies against representative
    data-leakage, data-corruption, and \rev{pointer redirection} violations
    (Section~\ref{subsec:eval_security_analysis}).
\end{itemize}

\gtan{Update the numbers below once the evaluation section is finalized.}
We evaluate \sysName{} on 10 library isolation boundaries, six in the
applications FFmpeg, memcached, Nginx, and rsync and four in programs that use
libpng, libucl, libyaml, and libzip (Section~\ref{sec:eval}). \rev{Averaged over the boundaries,
the inferred policies provide confidentiality coverage for 49.0\% of interface
fields and integrity coverage for 59.3\%. On average, 99.98\% of
pointer fields either need no runtime validation or have an inferred expected
domain (Section~\ref{subsec:eval_policy_coverage}). Relative to an RLBox sandboxing
baseline, \sysName{} adds 0.8\% on average to complete programs, and its
overhead is within measurement variation in 11 of 15 workloads. In library
stress workloads, which call the sandboxed library without application work
in between, it ranges from $-$6.7\%, where the policies remove a copy, to
8.6\% for queries that take about 120\,ns and return two pointers to check
(Section~\ref{subsec:eval_perf_overhead}).}
